\documentclass[11pt]{article}

\usepackage{nmrdoc}
\usepackage{siunitx}
\usepackage{bm}

\sisetup{per-mode=symbol}
\DeclareSIUnit{\angstrom}{\text{\AA}}
\emergencystretch=12em

\begin{document}
\NmrTitle{The WaterField Calculator}

\section{What the calculator computes}

The WaterField calculator describes the electrostatic field at each protein
atom due to explicit solvent water molecules in the same trajectory frame. It
uses the water oxygen and hydrogen coordinates, together with the point charges
assigned by the water model in the GROMACS topology, and accumulates two local
quantities at the target nucleus: a vector with units
\(\si{\volt\per\angstrom}\) and a rank-2 tensor with units
\(\si{\volt\per\square\angstrom}\).

The physical connection to shielding is the Buckingham electric-field effect.
The solvent electric environment perturbs the electronic distribution around a
nucleus. Because shielding is the electronic linear response to an applied
magnetic field, changing the electronic distribution can change the shielding
tensor. The vector field can feed scalar, element-dependent response terms in a
calibrated model, while the electric-field gradient carries a five-component
anisotropy that can be compared with the anisotropic part of DFT shielding.

The code does not apply Buckingham response coefficients and does not output a
chemical shift. Before calibration, the values are field-derived geometric
descriptors of the explicit-water configuration. The observable they bear on is
the nuclear shielding tensor, and through reference subtraction and sign
convention, the chemical shift.

\section{Where shielding enters}

For a nucleus in an applied magnetic field \(\bm B_0\), the local magnetic field
is written
\begin{equation}
  \bm B_{\mathrm{nuc}}=(\bm I-\bm\sigma)\bm B_0 .
  \label{eq:shielding}
\end{equation}
\(\bm I\) is the \(3\times3\) identity matrix. The shielding tensor
\(\bm\sigma\) is dimensionless: it is the proportionality between the applied
field and the secondary field produced at the nucleus by the electrons. A
matrix is needed because the response depends on orientation. A field applied
along one laboratory direction can induce a secondary field with components
along several molecular directions.

This is a quantum-mechanical response, not a classical electrostatics problem:
the applied magnetic field perturbs the electronic state, and the perturbed
state produces the magnetic field represented by \(\bm\sigma\). WaterField
does not solve that electronic problem. It computes local electrostatic
quantities that enter the usual field-response expansion of shielding,
\begin{equation}
  \sigma_{\mathrm{iso}}
    \approx \sigma_{\mathrm{iso}}^{(0)}
      + A_\alpha E_\alpha
      + B_{\alpha\beta}E_\alpha E_\beta+\cdots,
  \qquad
  \bm\sigma_{T2} \approx \Gamma[\bm V_{T2}] .
  \label{eq:buckingham-context}
\end{equation}
Repeated Cartesian indices are summed over \(x,y,z\). \(E_\alpha\) is a local
field component in \(\si{\volt\per\angstrom}\). \(\bm V_{T2}\) is the
traceless symmetric part of the electric-field-gradient tensor, in
\(\si{\volt\per\square\angstrom}\). The coefficients
\(A_\alpha\), \(B_{\alpha\beta}\), and the linear map \(\Gamma\) are not
implemented in this calculator. They stand for the calibrated, atom-type
dependent conversion from electrostatic descriptors to shielding in ppm.

The local picture behind the code is the shape of the electrostatic
potential. Let \(\phi(\bm x)\) be the potential in volts near a target atom at
\(\bm x_i\). For a small displacement \(\bm u\) from that atom,
\begin{equation}
  \phi(\bm x_i+\bm u)
   \approx \phi(\bm x_i)
      + \bm g_i\cdot\bm u
      + \frac{1}{2}\bm u^{\mathsf T}\bm V_i\bm u .
  \label{eq:taylor}
\end{equation}
\(\bm g_i=\nabla\phi(\bm x_i)\) is the slope vector of the potential surface,
and \(\bm V_i=\nabla\nabla\phi(\bm x_i)\) is its curvature matrix. This
second-order Taylor model says the vector term tells which direction the
potential rises fastest, and the matrix term tells how that slope changes with
direction. The code stores \(\bm g_i\) in the field channels. Textbook
electrostatics often defines the physical electric field as
\(\bm E_{\mathrm{phys}}=-\nabla\phi\), so the stored vector has the opposite
sign from that convention. The tensor \(\bm V_i\) is the conventional
potential Hessian used as the electric-field gradient in NMR.

\section{The tensor split used by the code}

A shielding tensor, an electric-field-gradient tensor, or any other
\(3\times3\) tensor \(\bm X\) is split by the C++ type
\texttt{SphericalTensor}. The scalar part is
\begin{equation}
  T_0=\frac{1}{3}\operatorname{tr}\bm X ,
  \label{eq:t0}
\end{equation}
the average of the diagonal entries. For shielding, this is the isotropic part
that survives molecular tumbling and sets the usual isotropic NMR shift after
referencing.

The antisymmetric part is stored as
\begin{equation}
  \bm T_1=\frac{1}{2}
  \begin{pmatrix}
    X_{yz}-X_{zy}\\
    X_{zx}-X_{xz}\\
    X_{xy}-X_{yx}
  \end{pmatrix}.
  \label{eq:t1}
\end{equation}
This three-vector is the dual of the part of the matrix that changes sign on
transpose. It is normally not observed directly in standard isotropic solution
NMR.

The remaining part is the symmetric traceless matrix
\begin{equation}
  \bm S=\frac{\bm X+\bm X^{\mathsf T}}{2}-T_0\bm I,
  \qquad \operatorname{tr}\bm S=0 .
  \label{eq:s}
\end{equation}
\(\bm S\) has five independent numbers. The code packs them as
\begin{equation}
  \bm T_2 =
  \begin{pmatrix}
    \sqrt{2}\,S_{xy}\\
    \sqrt{2}\,S_{yz}\\
    \sqrt{3/2}\,S_{zz}\\
    \sqrt{2}\,S_{xz}\\
    (S_{xx}-S_{yy})/\sqrt{2}
  \end{pmatrix}.
  \label{eq:t2}
\end{equation}
These factors are the code's norm-preserving real spherical normalization:
the squared length of the five-vector equals
\(\sum_{ab}S_{ab}^2\), the Frobenius norm of the matrix.

WaterField computes a vector channel and an EFG tensor channel. The EFG tensor
is built from symmetric outer products and then explicitly made traceless, so
its \(T_0\) and \(T_1\) components are structural zeros. The output files
therefore write only the five \(T_2\) components for
\texttt{water\_efg.npy} and \texttt{water\_efg\_first.npy}.

\section{The method implemented}

For a target protein atom \(i\) at position \(\bm x_i\), the solvent input is a
list of water molecules \(w\). Each water has three charge sites
\[
  (\bm s_{wO},q_{wO}),\qquad
  (\bm s_{wH1},q_{wH}),\qquad
  (\bm s_{wH2},q_{wH}),
\]
with positions in \(\si{\angstrom}\) and charges in elementary charge units.
The parameter file describes these as TIP3P water charge sites. For a standard
TIP3P topology, \(q_O=-0.834e\) and \(q_H=+0.417e\). The reader, however, does
not hard-code those numbers. It accepts a three-atom, one-O/two-H water
moltype in the GROMACS topology; the extraction path then uses the usual
O--H--H atom order and copies those topology charges into every
\texttt{WaterMolecule}. The actual charges in a run are therefore the charges
in the supplied topology.

For each target atom, the code scans the water list. A water is considered when
its oxygen lies within
\begin{equation}
  |\bm s_{wO}-\bm x_i| \le r_{\mathrm{cut}},
  \qquad r_{\mathrm{cut}}=\SI{15.0}{\angstrom}.
  \label{eq:oxygen-cutoff}
\end{equation}
This oxygen-distance test selects the molecule; once selected, all three charge
sites are evaluated. The first hydration shell is defined by
\(|\bm s_{wO}-\bm x_i|<\SI{3.5}{\angstrom}\). The second-shell count includes
waters outside the first shell but with oxygen distance below
\(\SI{5.5}{\angstrom}\). The second shell is a count, not a separate field
sum.

For one charge site \(j\) on a selected water, the implementation defines
\begin{equation}
  \bm r_{ij}=\bm s_j-\bm x_i,\qquad
  r_{ij}=|\bm r_{ij}|.
  \label{eq:r-def}
\end{equation}
\(\bm r_{ij}\) points from the target atom to the water charge site. If
\(r_{ij}<\SI{0.1}{\angstrom}\), the site is rejected by the shared
singularity guard because the Coulomb kernels diverge at zero distance.
Accepted sites contribute
\begin{equation}
  \bm g_i \leftarrow \bm g_i
  + k_e q_j\frac{\bm r_{ij}}{r_{ij}^3},
  \label{eq:g-sum}
\end{equation}
and
\begin{equation}
  \bm V_i \leftarrow \bm V_i
  + k_e q_j
    \left(
      \frac{3\bm r_{ij}\bm r_{ij}^{\mathsf T}}{r_{ij}^5}
      -\frac{\bm I}{r_{ij}^3}
    \right).
  \label{eq:v-sum}
\end{equation}
Here \(k_e=\SI{14.3996}{\volt\angstrom}\), \(q_j\) is in units of \(e\), and
the distances are in \(\si{\angstrom}\). Equation~\eqref{eq:g-sum} therefore
has units \(\si{\volt\per\angstrom}\), and Eq.~\eqref{eq:v-sum} has units
\(\si{\volt\per\square\angstrom}\). The dyadic product
\(\bm r\bm r^{\mathsf T}\) is the matrix with entries \(r_a r_b\).
Because the vector uses \(\bm s_j-\bm x_i\), its sign is opposite to the
usual textbook electric field of a positive point charge and opposite to the
protein \texttt{CoulombResult} vector convention. The tensor is unaffected by
that reversal because it depends on \(r_a r_b\).

Equation~\eqref{eq:v-sum} is the Hessian of the point-charge potential. Each
term is traceless away from the charge site: a Coulomb potential has zero
Laplacian where no source charge sits, and the trace of the Hessian is that
Laplacian. Floating-point summation leaves small trace residue, so the code
projects the totals back to the traceless subspace,
\begin{equation}
  \bm V_i \leftarrow
  \bm V_i-\frac{\operatorname{tr}\bm V_i}{3}\bm I,
  \qquad
  \bm V_i^{(1)} \leftarrow
  \bm V_i^{(1)}-\frac{\operatorname{tr}\bm V_i^{(1)}}{3}\bm I .
  \label{eq:detrace}
\end{equation}
\(\bm V_i^{(1)}\) is the same EFG sum restricted to waters whose oxygen is in
the first shell. Both \(\bm V_i\) and \(\bm V_i^{(1)}\) are then decomposed by
Eq.~\eqref{eq:t0}--\eqref{eq:t2}.

The vector channel has an additional numerical clamp. If
\begin{equation}
  |\bm g_i|>\SI{100.0}{\volt\per\angstrom},
  \label{eq:clamp-condition}
\end{equation}
then the code multiplies both the total vector \(\bm g_i\) and the first-shell
vector \(\bm g_i^{(1)}\) by
\begin{equation}
  c_i=\frac{\SI{100.0}{\volt\per\angstrom}}{|\bm g_i|}.
  \label{eq:clamp-scale}
\end{equation}
The EFG tensors are not scaled by this clamp. The clamp is a numerical sanity
guard for near-contact field magnitudes, not a physical saturation model.

The resulting feature files are \texttt{water\_efield.npy}, an
\((N,3)\) array containing the total vector \(\bm g_i\);
\texttt{water\_efield\_first.npy}, the first-shell vector;
\texttt{water\_efg.npy}, an \((N,5)\) array containing the total \(T_2\)
components of \(\bm V_i\); \texttt{water\_efg\_first.npy}, the first-shell
\(T_2\) components; and \texttt{water\_shell\_counts.npy}, an \((N,2)\) array
with first- and second-shell water oxygen counts.

The approximations follow from this code path. Waters are represented as fixed
point charges; no induced polarization, dielectric screening factor, or water
electronic response is evaluated inside the calculator. The support is a
finite oxygen-based cutoff sphere rather than an Ewald or PME electrostatic
sum. The water molecule is accepted or rejected by oxygen distance, although
the hydrogen sites may lie slightly closer or farther from the target. The
calculator scans explicit solvent coordinates for the current frame, so it
captures water placement and orientation in that frame, but it does not infer a
continuum solvent field.

\section{Inputs and output}

WaterField requires a \texttt{SolventEnvironment}. In trajectory mode that
environment is built from a full-system GROMACS TPR and coordinate frame. The
TPR supplies the water model topology and the oxygen and hydrogen point
charges; the coordinate frame supplies the explicit water positions, converted
from nanometres to \(\si{\angstrom}\). The reader recognizes three-atom,
one-O/two-H water moltypes by composition and then extracts positions in the
usual O--H--H order; four-site water models are not handled by this code path.

The calculator does not consume MOPAC charges, AIMNet2 charges or embeddings,
ff14SB protein charges, APBS continuum fields, or DSSP secondary-structure
annotations. Those data feed other calculators or downstream comparisons.
For WaterField, the external data are the explicit solvent coordinates and the
water point charges.
Ions may be present in the same \texttt{SolventEnvironment}, but this
calculator does not scan them.

The output is not a shielding tensor in ppm. It is a set of field-derived
geometric quantities: a vector slope of the water electrostatic potential, a
traceless rank-2 curvature tensor emitted through its five \(T_2\) components,
and hydration-shell occupancy counts. Calibration against DFT shielding
references is what assigns ppm-scale coefficients. Before that calibration, the
numbers should be read as explicit-solvent electrostatic descriptors whose
anisotropic tensor part can correlate with the anisotropic shielding response.

\end{document}
